% =====================================================================
%  significance.tex  --  PHASE V, Tasks 11 and 12
%
%  Significance testing and repeated runs: what the resampling unit is,
%  which differences survive, which designs cannot detect anything, and
%  what ten labelling seeds do to the headline state-time total.
%
%  Intended placement: Section 7 (\label{sec:significance}), AFTER both
%  result sections it tests -- the forecasting comparison and the
%  decision layer -- so that a reader has seen both tables before being
%  shown what survives resampling.  The decision layer's own intervals
%  are reported here rather than in decision_layer.tex, because the
%  resampling design they depend on is stated here.
%
%  Drop-in usage: % =====================================================================
%  significance.tex  --  PHASE V, Tasks 11 and 12
%
%  Significance testing and repeated runs: what the resampling unit is,
%  which differences survive, which designs cannot detect anything, and
%  what ten labelling seeds do to the headline state-time total.
%
%  Intended placement: Section 7 (\label{sec:significance}), AFTER both
%  result sections it tests -- the forecasting comparison and the
%  decision layer -- so that a reader has seen both tables before being
%  shown what survives resampling.  The decision layer's own intervals
%  are reported here rather than in decision_layer.tex, because the
%  resampling design they depend on is stated here.
%
%  Drop-in usage: % =====================================================================
%  significance.tex  --  PHASE V, Tasks 11 and 12
%
%  Significance testing and repeated runs: what the resampling unit is,
%  which differences survive, which designs cannot detect anything, and
%  what ten labelling seeds do to the headline state-time total.
%
%  Intended placement: Section 7 (\label{sec:significance}), AFTER both
%  result sections it tests -- the forecasting comparison and the
%  decision layer -- so that a reader has seen both tables before being
%  shown what survives resampling.  The decision layer's own intervals
%  are reported here rather than in decision_layer.tex, because the
%  resampling design they depend on is stated here.
%
%  Drop-in usage: % =====================================================================
%  significance.tex  --  PHASE V, Tasks 11 and 12
%
%  Significance testing and repeated runs: what the resampling unit is,
%  which differences survive, which designs cannot detect anything, and
%  what ten labelling seeds do to the headline state-time total.
%
%  Intended placement: Section 7 (\label{sec:significance}), AFTER both
%  result sections it tests -- the forecasting comparison and the
%  decision layer -- so that a reader has seen both tables before being
%  shown what survives resampling.  The decision layer's own intervals
%  are reported here rather than in decision_layer.tex, because the
%  resampling design they depend on is stated here.
%
%  Drop-in usage: \input{significance}
%  Figures are referenced by bare name; add to the preamble
%      \graphicspath{{../outputs/phase6/figures/}}
%  Bibliography: paper/references.bib
%  Numbers below are generated by  python run_phase5.py  and stored in
%    outputs/phase5/task11/pelletizer-I/forecasting_tests.json
%    outputs/phase5/task11/pelletizer-I/decision_bootstrap_S*.json
%    outputs/phase5/task11/cross_machine_tests.json
%    outputs/phase5/task11/state_block_bootstrap.csv
%    outputs/phase5/task12/pelletizer-I/{forecasting,labelling,decision}/
% =====================================================================

\section{Statistical Analysis}
\label{sec:significance}

Every headline quantity in the preceding sections is a point estimate.
This section reports what happens to those estimates under resampling
and under repeated fitting, and --- for two of the designs used here ---
what they are incapable of detecting in the first place.

\subsection{Four decisions that determine every interval below}
\label{sec:significance:method}

\textbf{The resampling unit is never a row.} A $1$\,Hz industrial record
is massively autocorrelated, and an i.i.d.\ bootstrap over rows would
treat $5.5$ million readings as $5.5$ million independent observations.
Measured on a synthetic AR(1) series with $\rho=0.999$, the i.i.d.\
bootstrap returns a standard error $25\times$ too small. Three
dependence-aware units are used instead, each chosen from the structure
of the quantity being estimated:

\begin{center}
\small
\begin{tabular}{@{}lll@{}}
\toprule
quantity & unit & why \\
\midrule
state-time totals & one-week moving blocks~\cite{kunsch1989jackknife}
  & duty cycle closes over a week \\
decision savings & factory days
  & the restart cap is per calendar day \\
forecast errors & test windows, \emph{paired}
  & one problem solved twice \\
\bottomrule
\end{tabular}
\end{center}

\textbf{A day drawn twice counts as two days.} In the decision bootstrap
the resampled copies are re-keyed. Leaving the label unchanged would let
two copies of one day share a single day's restart budget, silently
tightening the very constraint under evaluation.

\textbf{Every policy is scored on the same resample.} Drawing
independent resamples per policy would break the pairing and inflate the
interval of every contrast --- and the contrasts, not the individual
intervals, are what the decision claim rests on
(Section~\ref{sec:significance:decision}).

\textbf{Statistical and economic significance are separate columns.} A
paired test over $5\,449$ windows detects differences far below anything
that changes a decision, so every comparison carries a $p$-value, an
effect size, and a threshold fixed before the tests were run: $1$\,s of
per-window MAE, $0.02$ of $F_1$, $5$\,USD/yr. A result is allowed to be
\emph{detectable but immaterial}, and two of the twenty-one forecasting
contrasts are. Multiplicity is controlled by Holm's step-down
procedure~\cite{holm1979simple}, which achieves the same family-wise
error rate as Bonferroni and is uniformly more powerful;
Benjamini--Hochberg~\cite{benjamini1995controlling} is reported beside
it for the exploratory families, and both are shown next to the raw
$p$-value so the cost of the correction stays visible.

\subsection{Forecasting}
\label{sec:significance:forecast}

Seven models, $5\,449$ identical held-out windows. Every stored
prediction array was re-scored and matched against its recorded metrics
before any test ran, and all seven sit on a byte-identical vector of
targets, which is what makes the pairing legitimate.

\paragraph{The pre-registered test.}
The comparison named in advance was the proposed Seq2Seq LSTM against
the untrained reference. It resolves decisively, and in the direction
opposite to the one it was written to check
(Table~\ref{tab:sig:headline}). The design was set up to ask whether a
small \emph{positive} gap was real; there is no positive gap. The $F_1$
interval is computed on seed~$0$ of each model, which is the
\emph{best} of the sequence model's five seeds --- and even granting it
its most favourable draw the interval excludes zero in persistence's
favour. On the seed mean the gap is five times wider.

\begin{table}[t]
\centering
\caption{The pre-registered comparison: proposed sequence model versus
the untrained reference, on $5\,449$ paired windows.}
\label{tab:sig:headline}
\small
\begin{tabular}{@{}lrrr@{}}
\toprule
quantity & Seq2Seq & \textsc{persistence} & $p$ \\
\midrule
per-window MAE (s)        & $20.90$ & $12.95$ & $5.7\times10^{-44}$ \\
horizon-step accuracy     & \multicolumn{2}{c}{$-0.0574$ (worse)} & $1.5\times10^{-113}$ \\
STANDBY $F_1$, seed mean  & $0.561$ & $0.681$ & --- \\
STANDBY $F_1$, seed 0     & $0.660$ & $0.681$ & CI $[-0.042,-0.003]$ \\
\bottomrule
\end{tabular}
\end{table}

\paragraph{The family.}
Of the $21$ pairwise contrasts, $19$ are significant after Holm
correction and $17$ of those also clear the $1$\,s practical threshold.
Two conclusions survive both filters. \emph{Persistence beats all five
neural models}, at $p_{\text{Holm}}\le6\times10^{-31}$ and by $4.5$--$8.0$\,s of
MAE. \emph{XGBoost beats persistence}, by $1.45$\,s at
$p_{\text{Holm}}=7.0\times10^{-4}$ with rank-biserial correlation
$0.24$ --- the single contrast in the entire family in which training
beats not training.

\paragraph{Friedman over seeds, and how little it can see.}
Blocking on the five seeds and treating the six trained models as
treatments gives $\chi^2=12.77$, $p=0.026$, Kendall's $W=0.511$
\cite{friedman1937use,demsar2006statistical}. XGBoost takes mean rank
$1.00$ --- it is the best model on every seed block --- and Seq2Seq
takes $5.00$, the worst. But the Nemenyi critical difference is $3.37$
on a $1$--$6$ rank scale, so \textbf{the only pair this diagram
separates is XGBoost against Seq2Seq} ($4.00$ apart) and everything else
is joined (Figure~\ref{fig:cd}). Measured power at this design is $0.14$
for a $1$-sd effect and $0.53$ at $2$\,sd. The seed-blocked comparison
is nearly blind, and the confirmatory conclusions above do not rest on
it: those are paired over $5\,449$ windows, not over $5$ seeds.

\begin{figure}[t]
  \centering
  \includegraphics[width=\linewidth]{fig_p6_cd_diagram}
  \caption{Critical-difference diagram over five seed blocks. Models
  joined by a bar are not separable at $\alpha=0.05$. The bar spans
  almost the whole scale because $n=5$ blocks buys a critical difference
  of $3.37$ ranks.}
  \label{fig:cd}
\end{figure}

\subsection{The decision layer}
\label{sec:significance:decision}

Day-block bootstrap, $2\,000$ draws, on the minimum-dwell labelling:
$394$ held-out episodes over \textbf{18 factory days} spanning $73.6$
calendar days.

\begin{table}[t]
\centering
\caption{Annual saving against the plant's own behaviour (USD/yr per
machine), with $95\%$ day-block bootstrap intervals. Every policy is
scored on the same resample.}
\label{tab:sig:decision}
\small
\setlength{\tabcolsep}{3pt}
\begin{tabular}{@{}lrrr@{}}
\toprule
Policy & S1 ($10$\,min) & S2 ($60$\,min) & S3 ($240$\,min) \\
\midrule
never shut down   & $-219.8$\,{\footnotesize$[-334,-132]$} & $-157.1$\,{\footnotesize$[-276,-65]$} & $+68.6$\,{\footnotesize$[-84,+213]$} \\
immediate         & $-219.7$\,{\footnotesize$[-338,-126]$} & $-307.5$\,{\footnotesize$[-430,-212]$} & $-623.4$\,{\footnotesize$[-774,-478]$} \\
static break-even & $-166.7$\,{\footnotesize$[-290,-71]$}  & $-66.0$\,{\footnotesize$[-196,+24]$}  & $+6.3$\,{\footnotesize$[-170,+168]$} \\
ski-rental        & $-46.7$\,{\footnotesize$[-93,-7]$}     & $+0.0$\,{\footnotesize$[-25,+29]$}    & $+30.4$\,{\footnotesize$[-79,+155]$} \\
\textbf{proposed} & $\mathbf{-28.5}$\,{\footnotesize$[-79,+1]$} & $\mathbf{+2.9}$\,{\footnotesize$[-56,+46]$} & $\mathbf{+94.5}$\,{\footnotesize$[-58,+237]$} \\
oracle            & $+25.5$\,{\footnotesize$[+15,+38]$}    & $+37.6$\,{\footnotesize$[+17,+61]$}   & $+160.8$\,{\footnotesize$[+65,+266]$} \\
\bottomrule
\end{tabular}
\end{table}

Table~\ref{tab:sig:decision} has to be read twice, and the two readings
give different answers.

\textbf{Read row by row, no policy but the oracle demonstrates a
saving.} Every other interval against the status quo spans zero,
including the proposed policy's headline $+2.9$ USD/yr inside
$[-55.9,+46.2]$. The point estimate the earlier sections quote is
indistinguishable from zero, and always was; the bootstrap supplies the
interval that makes it visible.

\textbf{Read as paired contrasts, one result survives}
(Table~\ref{tab:sig:contrasts}). Because every policy is scored on the
same resample, a difference between two rows is far better determined
than either row: the pairing removes the day-to-day variance that makes
the individual intervals wide. The proposed policy beats the static
break-even rule it was built to replace by $+68.9$ USD/yr
($p=0.001$, interval excluding zero, above the $5$ USD/yr economic
threshold) in S2 and by $+138$ and $+88$ USD/yr in S1 and S3. It does
\emph{not} beat forecast-free ski-rental, and it cannot be shown to beat
the plant's own behaviour.

\begin{table}[t]
\centering
\caption{Paired contrasts, same resample. The claim the decision layer
supports is the third row, and only the third row.}
\label{tab:sig:contrasts}
\small
\begin{tabular}{@{}lrrr@{}}
\toprule
Contrast & S1 & S2 & S3 \\
\midrule
proposed $-$ status quo   & $-28.5$ & $+2.9$ & $+94.5$ \\
                          & \footnotesize$p=0.072$ & \footnotesize$p=0.828$ & \footnotesize$p=0.196$ \\
proposed $-$ ski-rental   & $+18.2$ & $+2.8$ & $+64.1$ \\
                          & \footnotesize$p=0.553$ & \footnotesize$p=0.729$ & \footnotesize$p=0.212$ \\
\textbf{proposed $-$ static} & $\mathbf{+138.2}$ & $\mathbf{+68.9}$ & $\mathbf{+88.3}$ \\
                          & \footnotesize$\mathbf{p<0.001}$ & \footnotesize$\mathbf{p=0.001}$ & \footnotesize$\mathbf{p<0.001}$ \\
proposed $-$ oracle       & $-54.0$ & $-34.8$ & $-66.3$ \\
                          & \footnotesize$p<0.001$ & \footnotesize$p<0.001$ & \footnotesize$p<0.001$ \\
\bottomrule
\end{tabular}
\end{table}

\textbf{The oracle gap is also real}, in all three scenarios: the
head-room is genuine and mostly uncaptured --- the proposed policy takes
$7.6\%$ of it in S2.

\textbf{The width is the sample, not the estimator.} Eighteen factory
days is the binding constraint: the day-block bootstrap has eighteen
units to resample, and no refinement of the estimator narrows these
intervals. Any claim of economic benefit from this layer needs a longer
record, and this paper does not make one.
Figure~\ref{fig:decision-forest} shows all three scenarios at once,
with the intervals that cross zero drawn so that they cannot be
mistaken for the ones that do not.

\begin{figure}[t]
  \centering
  \includegraphics[width=\linewidth]{fig_p6_decision_forest}
  \caption{Decision-layer intervals in the three scenarios. Intervals
  crossing zero are drawn grey and hollow; only the oracle's excludes
  zero in every scenario.}
  \label{fig:decision-forest}
\end{figure}

\subsection{Cross-machine tests, and a design that cannot detect
anything}
\label{sec:significance:cross}

Three paired comparisons across the eight machines are reported in
Table~\ref{tab:sig:cross}, and the third column of that table is as
important as the second.

\begin{table}[t]
\centering
\caption{Paired signed-rank tests across machines. The $p$-value floor
is the smallest two-sided value the test can return at that $n$.}
\label{tab:sig:cross}
\small
\setlength{\tabcolsep}{4pt}
\begin{tabular}{@{}lrrrrr@{}}
\toprule
Comparison & $n$ & $\Delta F_1$ & $p$ & floor & power \\
\midrule
own vs \textsc{persistence}, all      & 8 & $+0.194$ & $0.016$ & $0.008$ & $0.41$ \\
own vs \textsc{persistence}, motors   & 4 & $+0.048$ & $0.250$ & $\mathbf{0.125}$ & $\mathbf{0.00}$ \\
local vs transferred                  & 8 & $+0.384$ & $\mathbf{0.008}$ & $0.008$ & $d=1.75$ \\
\bottomrule
\end{tabular}
\end{table}

\textbf{Training beats not training across the eight machines, but the
effect lives in the auxiliary plant.} The $+0.194$ mean difference is
dominated by exhaust-fan~I ($0.982$ against $0.288$) and dpc~I ($0.402$
against $0.012$), machines whose ``STANDBY'' cluster is really their off
state and is therefore trivially predictable once a model has seen it
(Section~\ref{sec:multimachine}). On the four motor-driven machines, the
ones on which the state model is physically valid, the difference is
$+0.048$ and cannot be tested.

\textbf{``Cannot be tested'' is meant literally.} With four paired
observations the smallest two-sided $p$-value the signed-rank test can
return is $2/2^{4}=0.125$, which is above $\alpha=0.05$. No data, and no
effect size however large, can produce a significant result at $n=4$.
\textbf{A non-significant verdict here is a statement about the design,
not about the models}, and it is reported in exactly those terms rather
than as ``no significant difference''.
Figure~\ref{fig:power} puts numbers on both designs: at $n=8$ the paired
test reaches $80\%$ power only above $\approx1.4$\,sd; at $n=4$ it never
does, at any effect size.

\textbf{The transfer degradation is real.} Locally fitted sequence
models beat transferred ones by $0.384$ of $F_1$ with $p$ at the floor
and $d=+1.75$. This is the one cross-machine claim of
Section~\ref{sec:crossmachine} that survives testing outright.

\begin{figure}[t]
  \centering
  \includegraphics[width=\linewidth]{fig_p6_power_curve}
  \caption{Power of the paired signed-rank test against effect size for
  $n=4$ and $n=8$. The marker is the effect actually observed on the
  four motor-driven machines ($d=0.56$), whose measured power is $0.00$
  --- readable off the axis rather than taken on trust.}
  \label{fig:power}
\end{figure}

\subsection{Repeated fitting, and the result that matters most}
\label{sec:significance:seeds}

\subsubsection{The labelling is not reproducible from the seed alone}

Re-fitting the state model at ten seeds on one deterministic
preprocessing of pelletizer~I gives a STANDBY total of
$84.6\pm28.5$\,h --- a coefficient of variation of $\mathbf{33.7\%}$ on
the headline state-time quantity of the paper. Three seeds in ten
converge on a cluster that is not standby at all:

\begin{center}
\small
\begin{tabular}{@{}rrrrl@{}}
\toprule
seed & physics & STANDBY & mean power & what it found \\
\midrule
6 & $0.80$ & $121.6$\,h & $12\,961$\,W & a loaded running state \\
8 & $0.60$ & $63.2$\,h  & $0.8$\,W     & the OFF state \\
9 & $0.60$ & $13.8$\,h  & $1.1$\,W     & the OFF state \\
\bottomrule
\end{tabular}
\end{center}

A ``standby'' cluster at $0.8$\,W drawing $0.06\%$ of rated current is
the machine switched off; one at $13$\,kW is the machine working. EM
initialisation decides which, and in $30\%$ of initialisations it
decides wrongly.

\textbf{What rescues the result is the physics battery, and this is the
constructive finding.} The tests of
Section~\ref{sec:method:accept_reject} score exactly those three seeds
at $0.60$--$0.80$ and the other seven at $1.00$, so the accept/reject
decision is exact on this record. Conditioning on a passing score
collapses the spread:

\begin{center}
\small
\begin{tabular}{@{}lrr@{}}
\toprule
 & all $10$ seeds & the $7$ that pass \\
\midrule
STANDBY hours   & $84.6\pm28.5$\,h & $\mathbf{92.4\pm2.8}$\,h \\
coefficient of variation & $\mathbf{33.7\%}$ & $\mathbf{3.0\%}$ \\
range           & $13.8$--$121.6$\,h & $87.3$--$94.2$\,h \\
STANDBY power   & $4\,570\pm3\,543$\,W & $4\,676\pm343$\,W \\
\bottomrule
\end{tabular}
\end{center}

\textbf{The pipeline is reproducible because of the physics validation,
not despite the clustering.} That is a stronger statement of this
paper's contribution than an unconditional total would be, and it is why
the accept/reject loop is part of the method in
Section~\ref{sec:method:accept_reject} rather than an appendix to it.
Flicker is $0.00\%$ at every seed and the adjusted Rand index against
seed~$0$ stays $\ge0.889$: the \emph{partition} is stable because OFF
and PEAK\_LOAD dominate the record, and it is the identity of the
STANDBY cluster specifically that moves.

\subsubsection{Model-fitting variance is not what makes the economics
uncertain}

Ten seeds of the identical decision configuration (S2, minimum-dwell
labels) give a proposed-policy saving of
$\mathbf{+2.92\pm0.59}$\,USD/yr, $7.8\pm1.6\%$ of oracle, on head-room
of $37.63$\,USD/yr with \emph{no} seed variance at all --- head-room
being a property of the episodes and the coefficients rather than of the
fit. The forecaster's own error is $8\,954\pm156$\,s, about $2.5$\,h on
episode durations, which is why the policy captures under $8\%$ of what
is available.

Set the two sources of uncertainty side by side: model fitting
contributes $\pm0.6$\,USD/yr, day-to-day sampling contributes
$\pm50$\,USD/yr. \textbf{The uncertainty in the economic claim is the
18-day sample, not the model}, and no amount of retraining addresses it.

\subsubsection{State-time totals}

The one-week moving-block bootstrap gives intervals of $\pm15$--$25\%$
of the point estimate on the five-month records
(Table~\ref{tab:multimachine}, per machine). Combined with the $\pm3\%$
conditional labelling spread above, this fixes the precision at which
every state-time number in this paper is quoted: \textbf{two significant
figures, and conditional on the physics battery passing}. Three
significant figures assert a precision neither interval supports.

%  Figures are referenced by bare name; add to the preamble
%      \graphicspath{{../outputs/phase6/figures/}}
%  Bibliography: paper/references.bib
%  Numbers below are generated by  python run_phase5.py  and stored in
%    outputs/phase5/task11/pelletizer-I/forecasting_tests.json
%    outputs/phase5/task11/pelletizer-I/decision_bootstrap_S*.json
%    outputs/phase5/task11/cross_machine_tests.json
%    outputs/phase5/task11/state_block_bootstrap.csv
%    outputs/phase5/task12/pelletizer-I/{forecasting,labelling,decision}/
% =====================================================================

\section{Statistical Analysis}
\label{sec:significance}

Every headline quantity in the preceding sections is a point estimate.
This section reports what happens to those estimates under resampling
and under repeated fitting, and --- for two of the designs used here ---
what they are incapable of detecting in the first place.

\subsection{Four decisions that determine every interval below}
\label{sec:significance:method}

\textbf{The resampling unit is never a row.} A $1$\,Hz industrial record
is massively autocorrelated, and an i.i.d.\ bootstrap over rows would
treat $5.5$ million readings as $5.5$ million independent observations.
Measured on a synthetic AR(1) series with $\rho=0.999$, the i.i.d.\
bootstrap returns a standard error $25\times$ too small. Three
dependence-aware units are used instead, each chosen from the structure
of the quantity being estimated:

\begin{center}
\small
\begin{tabular}{@{}lll@{}}
\toprule
quantity & unit & why \\
\midrule
state-time totals & one-week moving blocks~\cite{kunsch1989jackknife}
  & duty cycle closes over a week \\
decision savings & factory days
  & the restart cap is per calendar day \\
forecast errors & test windows, \emph{paired}
  & one problem solved twice \\
\bottomrule
\end{tabular}
\end{center}

\textbf{A day drawn twice counts as two days.} In the decision bootstrap
the resampled copies are re-keyed. Leaving the label unchanged would let
two copies of one day share a single day's restart budget, silently
tightening the very constraint under evaluation.

\textbf{Every policy is scored on the same resample.} Drawing
independent resamples per policy would break the pairing and inflate the
interval of every contrast --- and the contrasts, not the individual
intervals, are what the decision claim rests on
(Section~\ref{sec:significance:decision}).

\textbf{Statistical and economic significance are separate columns.} A
paired test over $5\,449$ windows detects differences far below anything
that changes a decision, so every comparison carries a $p$-value, an
effect size, and a threshold fixed before the tests were run: $1$\,s of
per-window MAE, $0.02$ of $F_1$, $5$\,USD/yr. A result is allowed to be
\emph{detectable but immaterial}, and two of the twenty-one forecasting
contrasts are. Multiplicity is controlled by Holm's step-down
procedure~\cite{holm1979simple}, which achieves the same family-wise
error rate as Bonferroni and is uniformly more powerful;
Benjamini--Hochberg~\cite{benjamini1995controlling} is reported beside
it for the exploratory families, and both are shown next to the raw
$p$-value so the cost of the correction stays visible.

\subsection{Forecasting}
\label{sec:significance:forecast}

Seven models, $5\,449$ identical held-out windows. Every stored
prediction array was re-scored and matched against its recorded metrics
before any test ran, and all seven sit on a byte-identical vector of
targets, which is what makes the pairing legitimate.

\paragraph{The pre-registered test.}
The comparison named in advance was the proposed Seq2Seq LSTM against
the untrained reference. It resolves decisively, and in the direction
opposite to the one it was written to check
(Table~\ref{tab:sig:headline}). The design was set up to ask whether a
small \emph{positive} gap was real; there is no positive gap. The $F_1$
interval is computed on seed~$0$ of each model, which is the
\emph{best} of the sequence model's five seeds --- and even granting it
its most favourable draw the interval excludes zero in persistence's
favour. On the seed mean the gap is five times wider.

\begin{table}[t]
\centering
\caption{The pre-registered comparison: proposed sequence model versus
the untrained reference, on $5\,449$ paired windows.}
\label{tab:sig:headline}
\small
\begin{tabular}{@{}lrrr@{}}
\toprule
quantity & Seq2Seq & \textsc{persistence} & $p$ \\
\midrule
per-window MAE (s)        & $20.90$ & $12.95$ & $5.7\times10^{-44}$ \\
horizon-step accuracy     & \multicolumn{2}{c}{$-0.0574$ (worse)} & $1.5\times10^{-113}$ \\
STANDBY $F_1$, seed mean  & $0.561$ & $0.681$ & --- \\
STANDBY $F_1$, seed 0     & $0.660$ & $0.681$ & CI $[-0.042,-0.003]$ \\
\bottomrule
\end{tabular}
\end{table}

\paragraph{The family.}
Of the $21$ pairwise contrasts, $19$ are significant after Holm
correction and $17$ of those also clear the $1$\,s practical threshold.
Two conclusions survive both filters. \emph{Persistence beats all five
neural models}, at $p_{\text{Holm}}\le6\times10^{-31}$ and by $4.5$--$8.0$\,s of
MAE. \emph{XGBoost beats persistence}, by $1.45$\,s at
$p_{\text{Holm}}=7.0\times10^{-4}$ with rank-biserial correlation
$0.24$ --- the single contrast in the entire family in which training
beats not training.

\paragraph{Friedman over seeds, and how little it can see.}
Blocking on the five seeds and treating the six trained models as
treatments gives $\chi^2=12.77$, $p=0.026$, Kendall's $W=0.511$
\cite{friedman1937use,demsar2006statistical}. XGBoost takes mean rank
$1.00$ --- it is the best model on every seed block --- and Seq2Seq
takes $5.00$, the worst. But the Nemenyi critical difference is $3.37$
on a $1$--$6$ rank scale, so \textbf{the only pair this diagram
separates is XGBoost against Seq2Seq} ($4.00$ apart) and everything else
is joined (Figure~\ref{fig:cd}). Measured power at this design is $0.14$
for a $1$-sd effect and $0.53$ at $2$\,sd. The seed-blocked comparison
is nearly blind, and the confirmatory conclusions above do not rest on
it: those are paired over $5\,449$ windows, not over $5$ seeds.

\begin{figure}[t]
  \centering
  \includegraphics[width=\linewidth]{fig_p6_cd_diagram}
  \caption{Critical-difference diagram over five seed blocks. Models
  joined by a bar are not separable at $\alpha=0.05$. The bar spans
  almost the whole scale because $n=5$ blocks buys a critical difference
  of $3.37$ ranks.}
  \label{fig:cd}
\end{figure}

\subsection{The decision layer}
\label{sec:significance:decision}

Day-block bootstrap, $2\,000$ draws, on the minimum-dwell labelling:
$394$ held-out episodes over \textbf{18 factory days} spanning $73.6$
calendar days.

\begin{table}[t]
\centering
\caption{Annual saving against the plant's own behaviour (USD/yr per
machine), with $95\%$ day-block bootstrap intervals. Every policy is
scored on the same resample.}
\label{tab:sig:decision}
\small
\setlength{\tabcolsep}{3pt}
\begin{tabular}{@{}lrrr@{}}
\toprule
Policy & S1 ($10$\,min) & S2 ($60$\,min) & S3 ($240$\,min) \\
\midrule
never shut down   & $-219.8$\,{\footnotesize$[-334,-132]$} & $-157.1$\,{\footnotesize$[-276,-65]$} & $+68.6$\,{\footnotesize$[-84,+213]$} \\
immediate         & $-219.7$\,{\footnotesize$[-338,-126]$} & $-307.5$\,{\footnotesize$[-430,-212]$} & $-623.4$\,{\footnotesize$[-774,-478]$} \\
static break-even & $-166.7$\,{\footnotesize$[-290,-71]$}  & $-66.0$\,{\footnotesize$[-196,+24]$}  & $+6.3$\,{\footnotesize$[-170,+168]$} \\
ski-rental        & $-46.7$\,{\footnotesize$[-93,-7]$}     & $+0.0$\,{\footnotesize$[-25,+29]$}    & $+30.4$\,{\footnotesize$[-79,+155]$} \\
\textbf{proposed} & $\mathbf{-28.5}$\,{\footnotesize$[-79,+1]$} & $\mathbf{+2.9}$\,{\footnotesize$[-56,+46]$} & $\mathbf{+94.5}$\,{\footnotesize$[-58,+237]$} \\
oracle            & $+25.5$\,{\footnotesize$[+15,+38]$}    & $+37.6$\,{\footnotesize$[+17,+61]$}   & $+160.8$\,{\footnotesize$[+65,+266]$} \\
\bottomrule
\end{tabular}
\end{table}

Table~\ref{tab:sig:decision} has to be read twice, and the two readings
give different answers.

\textbf{Read row by row, no policy but the oracle demonstrates a
saving.} Every other interval against the status quo spans zero,
including the proposed policy's headline $+2.9$ USD/yr inside
$[-55.9,+46.2]$. The point estimate the earlier sections quote is
indistinguishable from zero, and always was; the bootstrap supplies the
interval that makes it visible.

\textbf{Read as paired contrasts, one result survives}
(Table~\ref{tab:sig:contrasts}). Because every policy is scored on the
same resample, a difference between two rows is far better determined
than either row: the pairing removes the day-to-day variance that makes
the individual intervals wide. The proposed policy beats the static
break-even rule it was built to replace by $+68.9$ USD/yr
($p=0.001$, interval excluding zero, above the $5$ USD/yr economic
threshold) in S2 and by $+138$ and $+88$ USD/yr in S1 and S3. It does
\emph{not} beat forecast-free ski-rental, and it cannot be shown to beat
the plant's own behaviour.

\begin{table}[t]
\centering
\caption{Paired contrasts, same resample. The claim the decision layer
supports is the third row, and only the third row.}
\label{tab:sig:contrasts}
\small
\begin{tabular}{@{}lrrr@{}}
\toprule
Contrast & S1 & S2 & S3 \\
\midrule
proposed $-$ status quo   & $-28.5$ & $+2.9$ & $+94.5$ \\
                          & \footnotesize$p=0.072$ & \footnotesize$p=0.828$ & \footnotesize$p=0.196$ \\
proposed $-$ ski-rental   & $+18.2$ & $+2.8$ & $+64.1$ \\
                          & \footnotesize$p=0.553$ & \footnotesize$p=0.729$ & \footnotesize$p=0.212$ \\
\textbf{proposed $-$ static} & $\mathbf{+138.2}$ & $\mathbf{+68.9}$ & $\mathbf{+88.3}$ \\
                          & \footnotesize$\mathbf{p<0.001}$ & \footnotesize$\mathbf{p=0.001}$ & \footnotesize$\mathbf{p<0.001}$ \\
proposed $-$ oracle       & $-54.0$ & $-34.8$ & $-66.3$ \\
                          & \footnotesize$p<0.001$ & \footnotesize$p<0.001$ & \footnotesize$p<0.001$ \\
\bottomrule
\end{tabular}
\end{table}

\textbf{The oracle gap is also real}, in all three scenarios: the
head-room is genuine and mostly uncaptured --- the proposed policy takes
$7.6\%$ of it in S2.

\textbf{The width is the sample, not the estimator.} Eighteen factory
days is the binding constraint: the day-block bootstrap has eighteen
units to resample, and no refinement of the estimator narrows these
intervals. Any claim of economic benefit from this layer needs a longer
record, and this paper does not make one.
Figure~\ref{fig:decision-forest} shows all three scenarios at once,
with the intervals that cross zero drawn so that they cannot be
mistaken for the ones that do not.

\begin{figure}[t]
  \centering
  \includegraphics[width=\linewidth]{fig_p6_decision_forest}
  \caption{Decision-layer intervals in the three scenarios. Intervals
  crossing zero are drawn grey and hollow; only the oracle's excludes
  zero in every scenario.}
  \label{fig:decision-forest}
\end{figure}

\subsection{Cross-machine tests, and a design that cannot detect
anything}
\label{sec:significance:cross}

Three paired comparisons across the eight machines are reported in
Table~\ref{tab:sig:cross}, and the third column of that table is as
important as the second.

\begin{table}[t]
\centering
\caption{Paired signed-rank tests across machines. The $p$-value floor
is the smallest two-sided value the test can return at that $n$.}
\label{tab:sig:cross}
\small
\setlength{\tabcolsep}{4pt}
\begin{tabular}{@{}lrrrrr@{}}
\toprule
Comparison & $n$ & $\Delta F_1$ & $p$ & floor & power \\
\midrule
own vs \textsc{persistence}, all      & 8 & $+0.194$ & $0.016$ & $0.008$ & $0.41$ \\
own vs \textsc{persistence}, motors   & 4 & $+0.048$ & $0.250$ & $\mathbf{0.125}$ & $\mathbf{0.00}$ \\
local vs transferred                  & 8 & $+0.384$ & $\mathbf{0.008}$ & $0.008$ & $d=1.75$ \\
\bottomrule
\end{tabular}
\end{table}

\textbf{Training beats not training across the eight machines, but the
effect lives in the auxiliary plant.} The $+0.194$ mean difference is
dominated by exhaust-fan~I ($0.982$ against $0.288$) and dpc~I ($0.402$
against $0.012$), machines whose ``STANDBY'' cluster is really their off
state and is therefore trivially predictable once a model has seen it
(Section~\ref{sec:multimachine}). On the four motor-driven machines, the
ones on which the state model is physically valid, the difference is
$+0.048$ and cannot be tested.

\textbf{``Cannot be tested'' is meant literally.} With four paired
observations the smallest two-sided $p$-value the signed-rank test can
return is $2/2^{4}=0.125$, which is above $\alpha=0.05$. No data, and no
effect size however large, can produce a significant result at $n=4$.
\textbf{A non-significant verdict here is a statement about the design,
not about the models}, and it is reported in exactly those terms rather
than as ``no significant difference''.
Figure~\ref{fig:power} puts numbers on both designs: at $n=8$ the paired
test reaches $80\%$ power only above $\approx1.4$\,sd; at $n=4$ it never
does, at any effect size.

\textbf{The transfer degradation is real.} Locally fitted sequence
models beat transferred ones by $0.384$ of $F_1$ with $p$ at the floor
and $d=+1.75$. This is the one cross-machine claim of
Section~\ref{sec:crossmachine} that survives testing outright.

\begin{figure}[t]
  \centering
  \includegraphics[width=\linewidth]{fig_p6_power_curve}
  \caption{Power of the paired signed-rank test against effect size for
  $n=4$ and $n=8$. The marker is the effect actually observed on the
  four motor-driven machines ($d=0.56$), whose measured power is $0.00$
  --- readable off the axis rather than taken on trust.}
  \label{fig:power}
\end{figure}

\subsection{Repeated fitting, and the result that matters most}
\label{sec:significance:seeds}

\subsubsection{The labelling is not reproducible from the seed alone}

Re-fitting the state model at ten seeds on one deterministic
preprocessing of pelletizer~I gives a STANDBY total of
$84.6\pm28.5$\,h --- a coefficient of variation of $\mathbf{33.7\%}$ on
the headline state-time quantity of the paper. Three seeds in ten
converge on a cluster that is not standby at all:

\begin{center}
\small
\begin{tabular}{@{}rrrrl@{}}
\toprule
seed & physics & STANDBY & mean power & what it found \\
\midrule
6 & $0.80$ & $121.6$\,h & $12\,961$\,W & a loaded running state \\
8 & $0.60$ & $63.2$\,h  & $0.8$\,W     & the OFF state \\
9 & $0.60$ & $13.8$\,h  & $1.1$\,W     & the OFF state \\
\bottomrule
\end{tabular}
\end{center}

A ``standby'' cluster at $0.8$\,W drawing $0.06\%$ of rated current is
the machine switched off; one at $13$\,kW is the machine working. EM
initialisation decides which, and in $30\%$ of initialisations it
decides wrongly.

\textbf{What rescues the result is the physics battery, and this is the
constructive finding.} The tests of
Section~\ref{sec:method:accept_reject} score exactly those three seeds
at $0.60$--$0.80$ and the other seven at $1.00$, so the accept/reject
decision is exact on this record. Conditioning on a passing score
collapses the spread:

\begin{center}
\small
\begin{tabular}{@{}lrr@{}}
\toprule
 & all $10$ seeds & the $7$ that pass \\
\midrule
STANDBY hours   & $84.6\pm28.5$\,h & $\mathbf{92.4\pm2.8}$\,h \\
coefficient of variation & $\mathbf{33.7\%}$ & $\mathbf{3.0\%}$ \\
range           & $13.8$--$121.6$\,h & $87.3$--$94.2$\,h \\
STANDBY power   & $4\,570\pm3\,543$\,W & $4\,676\pm343$\,W \\
\bottomrule
\end{tabular}
\end{center}

\textbf{The pipeline is reproducible because of the physics validation,
not despite the clustering.} That is a stronger statement of this
paper's contribution than an unconditional total would be, and it is why
the accept/reject loop is part of the method in
Section~\ref{sec:method:accept_reject} rather than an appendix to it.
Flicker is $0.00\%$ at every seed and the adjusted Rand index against
seed~$0$ stays $\ge0.889$: the \emph{partition} is stable because OFF
and PEAK\_LOAD dominate the record, and it is the identity of the
STANDBY cluster specifically that moves.

\subsubsection{Model-fitting variance is not what makes the economics
uncertain}

Ten seeds of the identical decision configuration (S2, minimum-dwell
labels) give a proposed-policy saving of
$\mathbf{+2.92\pm0.59}$\,USD/yr, $7.8\pm1.6\%$ of oracle, on head-room
of $37.63$\,USD/yr with \emph{no} seed variance at all --- head-room
being a property of the episodes and the coefficients rather than of the
fit. The forecaster's own error is $8\,954\pm156$\,s, about $2.5$\,h on
episode durations, which is why the policy captures under $8\%$ of what
is available.

Set the two sources of uncertainty side by side: model fitting
contributes $\pm0.6$\,USD/yr, day-to-day sampling contributes
$\pm50$\,USD/yr. \textbf{The uncertainty in the economic claim is the
18-day sample, not the model}, and no amount of retraining addresses it.

\subsubsection{State-time totals}

The one-week moving-block bootstrap gives intervals of $\pm15$--$25\%$
of the point estimate on the five-month records
(Table~\ref{tab:multimachine}, per machine). Combined with the $\pm3\%$
conditional labelling spread above, this fixes the precision at which
every state-time number in this paper is quoted: \textbf{two significant
figures, and conditional on the physics battery passing}. Three
significant figures assert a precision neither interval supports.

%  Figures are referenced by bare name; add to the preamble
%      \graphicspath{{../outputs/phase6/figures/}}
%  Bibliography: paper/references.bib
%  Numbers below are generated by  python run_phase5.py  and stored in
%    outputs/phase5/task11/pelletizer-I/forecasting_tests.json
%    outputs/phase5/task11/pelletizer-I/decision_bootstrap_S*.json
%    outputs/phase5/task11/cross_machine_tests.json
%    outputs/phase5/task11/state_block_bootstrap.csv
%    outputs/phase5/task12/pelletizer-I/{forecasting,labelling,decision}/
% =====================================================================

\section{Statistical Analysis}
\label{sec:significance}

Every headline quantity in the preceding sections is a point estimate.
This section reports what happens to those estimates under resampling
and under repeated fitting, and --- for two of the designs used here ---
what they are incapable of detecting in the first place.

\subsection{Four decisions that determine every interval below}
\label{sec:significance:method}

\textbf{The resampling unit is never a row.} A $1$\,Hz industrial record
is massively autocorrelated, and an i.i.d.\ bootstrap over rows would
treat $5.5$ million readings as $5.5$ million independent observations.
Measured on a synthetic AR(1) series with $\rho=0.999$, the i.i.d.\
bootstrap returns a standard error $25\times$ too small. Three
dependence-aware units are used instead, each chosen from the structure
of the quantity being estimated:

\begin{center}
\small
\begin{tabular}{@{}lll@{}}
\toprule
quantity & unit & why \\
\midrule
state-time totals & one-week moving blocks~\cite{kunsch1989jackknife}
  & duty cycle closes over a week \\
decision savings & factory days
  & the restart cap is per calendar day \\
forecast errors & test windows, \emph{paired}
  & one problem solved twice \\
\bottomrule
\end{tabular}
\end{center}

\textbf{A day drawn twice counts as two days.} In the decision bootstrap
the resampled copies are re-keyed. Leaving the label unchanged would let
two copies of one day share a single day's restart budget, silently
tightening the very constraint under evaluation.

\textbf{Every policy is scored on the same resample.} Drawing
independent resamples per policy would break the pairing and inflate the
interval of every contrast --- and the contrasts, not the individual
intervals, are what the decision claim rests on
(Section~\ref{sec:significance:decision}).

\textbf{Statistical and economic significance are separate columns.} A
paired test over $5\,449$ windows detects differences far below anything
that changes a decision, so every comparison carries a $p$-value, an
effect size, and a threshold fixed before the tests were run: $1$\,s of
per-window MAE, $0.02$ of $F_1$, $5$\,USD/yr. A result is allowed to be
\emph{detectable but immaterial}, and two of the twenty-one forecasting
contrasts are. Multiplicity is controlled by Holm's step-down
procedure~\cite{holm1979simple}, which achieves the same family-wise
error rate as Bonferroni and is uniformly more powerful;
Benjamini--Hochberg~\cite{benjamini1995controlling} is reported beside
it for the exploratory families, and both are shown next to the raw
$p$-value so the cost of the correction stays visible.

\subsection{Forecasting}
\label{sec:significance:forecast}

Seven models, $5\,449$ identical held-out windows. Every stored
prediction array was re-scored and matched against its recorded metrics
before any test ran, and all seven sit on a byte-identical vector of
targets, which is what makes the pairing legitimate.

\paragraph{The pre-registered test.}
The comparison named in advance was the proposed Seq2Seq LSTM against
the untrained reference. It resolves decisively, and in the direction
opposite to the one it was written to check
(Table~\ref{tab:sig:headline}). The design was set up to ask whether a
small \emph{positive} gap was real; there is no positive gap. The $F_1$
interval is computed on seed~$0$ of each model, which is the
\emph{best} of the sequence model's five seeds --- and even granting it
its most favourable draw the interval excludes zero in persistence's
favour. On the seed mean the gap is five times wider.

\begin{table}[t]
\centering
\caption{The pre-registered comparison: proposed sequence model versus
the untrained reference, on $5\,449$ paired windows.}
\label{tab:sig:headline}
\small
\begin{tabular}{@{}lrrr@{}}
\toprule
quantity & Seq2Seq & \textsc{persistence} & $p$ \\
\midrule
per-window MAE (s)        & $20.90$ & $12.95$ & $5.7\times10^{-44}$ \\
horizon-step accuracy     & \multicolumn{2}{c}{$-0.0574$ (worse)} & $1.5\times10^{-113}$ \\
STANDBY $F_1$, seed mean  & $0.561$ & $0.681$ & --- \\
STANDBY $F_1$, seed 0     & $0.660$ & $0.681$ & CI $[-0.042,-0.003]$ \\
\bottomrule
\end{tabular}
\end{table}

\paragraph{The family.}
Of the $21$ pairwise contrasts, $19$ are significant after Holm
correction and $17$ of those also clear the $1$\,s practical threshold.
Two conclusions survive both filters. \emph{Persistence beats all five
neural models}, at $p_{\text{Holm}}\le6\times10^{-31}$ and by $4.5$--$8.0$\,s of
MAE. \emph{XGBoost beats persistence}, by $1.45$\,s at
$p_{\text{Holm}}=7.0\times10^{-4}$ with rank-biserial correlation
$0.24$ --- the single contrast in the entire family in which training
beats not training.

\paragraph{Friedman over seeds, and how little it can see.}
Blocking on the five seeds and treating the six trained models as
treatments gives $\chi^2=12.77$, $p=0.026$, Kendall's $W=0.511$
\cite{friedman1937use,demsar2006statistical}. XGBoost takes mean rank
$1.00$ --- it is the best model on every seed block --- and Seq2Seq
takes $5.00$, the worst. But the Nemenyi critical difference is $3.37$
on a $1$--$6$ rank scale, so \textbf{the only pair this diagram
separates is XGBoost against Seq2Seq} ($4.00$ apart) and everything else
is joined (Figure~\ref{fig:cd}). Measured power at this design is $0.14$
for a $1$-sd effect and $0.53$ at $2$\,sd. The seed-blocked comparison
is nearly blind, and the confirmatory conclusions above do not rest on
it: those are paired over $5\,449$ windows, not over $5$ seeds.

\begin{figure}[t]
  \centering
  \includegraphics[width=\linewidth]{fig_p6_cd_diagram}
  \caption{Critical-difference diagram over five seed blocks. Models
  joined by a bar are not separable at $\alpha=0.05$. The bar spans
  almost the whole scale because $n=5$ blocks buys a critical difference
  of $3.37$ ranks.}
  \label{fig:cd}
\end{figure}

\subsection{The decision layer}
\label{sec:significance:decision}

Day-block bootstrap, $2\,000$ draws, on the minimum-dwell labelling:
$394$ held-out episodes over \textbf{18 factory days} spanning $73.6$
calendar days.

\begin{table}[t]
\centering
\caption{Annual saving against the plant's own behaviour (USD/yr per
machine), with $95\%$ day-block bootstrap intervals. Every policy is
scored on the same resample.}
\label{tab:sig:decision}
\small
\setlength{\tabcolsep}{3pt}
\begin{tabular}{@{}lrrr@{}}
\toprule
Policy & S1 ($10$\,min) & S2 ($60$\,min) & S3 ($240$\,min) \\
\midrule
never shut down   & $-219.8$\,{\footnotesize$[-334,-132]$} & $-157.1$\,{\footnotesize$[-276,-65]$} & $+68.6$\,{\footnotesize$[-84,+213]$} \\
immediate         & $-219.7$\,{\footnotesize$[-338,-126]$} & $-307.5$\,{\footnotesize$[-430,-212]$} & $-623.4$\,{\footnotesize$[-774,-478]$} \\
static break-even & $-166.7$\,{\footnotesize$[-290,-71]$}  & $-66.0$\,{\footnotesize$[-196,+24]$}  & $+6.3$\,{\footnotesize$[-170,+168]$} \\
ski-rental        & $-46.7$\,{\footnotesize$[-93,-7]$}     & $+0.0$\,{\footnotesize$[-25,+29]$}    & $+30.4$\,{\footnotesize$[-79,+155]$} \\
\textbf{proposed} & $\mathbf{-28.5}$\,{\footnotesize$[-79,+1]$} & $\mathbf{+2.9}$\,{\footnotesize$[-56,+46]$} & $\mathbf{+94.5}$\,{\footnotesize$[-58,+237]$} \\
oracle            & $+25.5$\,{\footnotesize$[+15,+38]$}    & $+37.6$\,{\footnotesize$[+17,+61]$}   & $+160.8$\,{\footnotesize$[+65,+266]$} \\
\bottomrule
\end{tabular}
\end{table}

Table~\ref{tab:sig:decision} has to be read twice, and the two readings
give different answers.

\textbf{Read row by row, no policy but the oracle demonstrates a
saving.} Every other interval against the status quo spans zero,
including the proposed policy's headline $+2.9$ USD/yr inside
$[-55.9,+46.2]$. The point estimate the earlier sections quote is
indistinguishable from zero, and always was; the bootstrap supplies the
interval that makes it visible.

\textbf{Read as paired contrasts, one result survives}
(Table~\ref{tab:sig:contrasts}). Because every policy is scored on the
same resample, a difference between two rows is far better determined
than either row: the pairing removes the day-to-day variance that makes
the individual intervals wide. The proposed policy beats the static
break-even rule it was built to replace by $+68.9$ USD/yr
($p=0.001$, interval excluding zero, above the $5$ USD/yr economic
threshold) in S2 and by $+138$ and $+88$ USD/yr in S1 and S3. It does
\emph{not} beat forecast-free ski-rental, and it cannot be shown to beat
the plant's own behaviour.

\begin{table}[t]
\centering
\caption{Paired contrasts, same resample. The claim the decision layer
supports is the third row, and only the third row.}
\label{tab:sig:contrasts}
\small
\begin{tabular}{@{}lrrr@{}}
\toprule
Contrast & S1 & S2 & S3 \\
\midrule
proposed $-$ status quo   & $-28.5$ & $+2.9$ & $+94.5$ \\
                          & \footnotesize$p=0.072$ & \footnotesize$p=0.828$ & \footnotesize$p=0.196$ \\
proposed $-$ ski-rental   & $+18.2$ & $+2.8$ & $+64.1$ \\
                          & \footnotesize$p=0.553$ & \footnotesize$p=0.729$ & \footnotesize$p=0.212$ \\
\textbf{proposed $-$ static} & $\mathbf{+138.2}$ & $\mathbf{+68.9}$ & $\mathbf{+88.3}$ \\
                          & \footnotesize$\mathbf{p<0.001}$ & \footnotesize$\mathbf{p=0.001}$ & \footnotesize$\mathbf{p<0.001}$ \\
proposed $-$ oracle       & $-54.0$ & $-34.8$ & $-66.3$ \\
                          & \footnotesize$p<0.001$ & \footnotesize$p<0.001$ & \footnotesize$p<0.001$ \\
\bottomrule
\end{tabular}
\end{table}

\textbf{The oracle gap is also real}, in all three scenarios: the
head-room is genuine and mostly uncaptured --- the proposed policy takes
$7.6\%$ of it in S2.

\textbf{The width is the sample, not the estimator.} Eighteen factory
days is the binding constraint: the day-block bootstrap has eighteen
units to resample, and no refinement of the estimator narrows these
intervals. Any claim of economic benefit from this layer needs a longer
record, and this paper does not make one.
Figure~\ref{fig:decision-forest} shows all three scenarios at once,
with the intervals that cross zero drawn so that they cannot be
mistaken for the ones that do not.

\begin{figure}[t]
  \centering
  \includegraphics[width=\linewidth]{fig_p6_decision_forest}
  \caption{Decision-layer intervals in the three scenarios. Intervals
  crossing zero are drawn grey and hollow; only the oracle's excludes
  zero in every scenario.}
  \label{fig:decision-forest}
\end{figure}

\subsection{Cross-machine tests, and a design that cannot detect
anything}
\label{sec:significance:cross}

Three paired comparisons across the eight machines are reported in
Table~\ref{tab:sig:cross}, and the third column of that table is as
important as the second.

\begin{table}[t]
\centering
\caption{Paired signed-rank tests across machines. The $p$-value floor
is the smallest two-sided value the test can return at that $n$.}
\label{tab:sig:cross}
\small
\setlength{\tabcolsep}{4pt}
\begin{tabular}{@{}lrrrrr@{}}
\toprule
Comparison & $n$ & $\Delta F_1$ & $p$ & floor & power \\
\midrule
own vs \textsc{persistence}, all      & 8 & $+0.194$ & $0.016$ & $0.008$ & $0.41$ \\
own vs \textsc{persistence}, motors   & 4 & $+0.048$ & $0.250$ & $\mathbf{0.125}$ & $\mathbf{0.00}$ \\
local vs transferred                  & 8 & $+0.384$ & $\mathbf{0.008}$ & $0.008$ & $d=1.75$ \\
\bottomrule
\end{tabular}
\end{table}

\textbf{Training beats not training across the eight machines, but the
effect lives in the auxiliary plant.} The $+0.194$ mean difference is
dominated by exhaust-fan~I ($0.982$ against $0.288$) and dpc~I ($0.402$
against $0.012$), machines whose ``STANDBY'' cluster is really their off
state and is therefore trivially predictable once a model has seen it
(Section~\ref{sec:multimachine}). On the four motor-driven machines, the
ones on which the state model is physically valid, the difference is
$+0.048$ and cannot be tested.

\textbf{``Cannot be tested'' is meant literally.} With four paired
observations the smallest two-sided $p$-value the signed-rank test can
return is $2/2^{4}=0.125$, which is above $\alpha=0.05$. No data, and no
effect size however large, can produce a significant result at $n=4$.
\textbf{A non-significant verdict here is a statement about the design,
not about the models}, and it is reported in exactly those terms rather
than as ``no significant difference''.
Figure~\ref{fig:power} puts numbers on both designs: at $n=8$ the paired
test reaches $80\%$ power only above $\approx1.4$\,sd; at $n=4$ it never
does, at any effect size.

\textbf{The transfer degradation is real.} Locally fitted sequence
models beat transferred ones by $0.384$ of $F_1$ with $p$ at the floor
and $d=+1.75$. This is the one cross-machine claim of
Section~\ref{sec:crossmachine} that survives testing outright.

\begin{figure}[t]
  \centering
  \includegraphics[width=\linewidth]{fig_p6_power_curve}
  \caption{Power of the paired signed-rank test against effect size for
  $n=4$ and $n=8$. The marker is the effect actually observed on the
  four motor-driven machines ($d=0.56$), whose measured power is $0.00$
  --- readable off the axis rather than taken on trust.}
  \label{fig:power}
\end{figure}

\subsection{Repeated fitting, and the result that matters most}
\label{sec:significance:seeds}

\subsubsection{The labelling is not reproducible from the seed alone}

Re-fitting the state model at ten seeds on one deterministic
preprocessing of pelletizer~I gives a STANDBY total of
$84.6\pm28.5$\,h --- a coefficient of variation of $\mathbf{33.7\%}$ on
the headline state-time quantity of the paper. Three seeds in ten
converge on a cluster that is not standby at all:

\begin{center}
\small
\begin{tabular}{@{}rrrrl@{}}
\toprule
seed & physics & STANDBY & mean power & what it found \\
\midrule
6 & $0.80$ & $121.6$\,h & $12\,961$\,W & a loaded running state \\
8 & $0.60$ & $63.2$\,h  & $0.8$\,W     & the OFF state \\
9 & $0.60$ & $13.8$\,h  & $1.1$\,W     & the OFF state \\
\bottomrule
\end{tabular}
\end{center}

A ``standby'' cluster at $0.8$\,W drawing $0.06\%$ of rated current is
the machine switched off; one at $13$\,kW is the machine working. EM
initialisation decides which, and in $30\%$ of initialisations it
decides wrongly.

\textbf{What rescues the result is the physics battery, and this is the
constructive finding.} The tests of
Section~\ref{sec:method:accept_reject} score exactly those three seeds
at $0.60$--$0.80$ and the other seven at $1.00$, so the accept/reject
decision is exact on this record. Conditioning on a passing score
collapses the spread:

\begin{center}
\small
\begin{tabular}{@{}lrr@{}}
\toprule
 & all $10$ seeds & the $7$ that pass \\
\midrule
STANDBY hours   & $84.6\pm28.5$\,h & $\mathbf{92.4\pm2.8}$\,h \\
coefficient of variation & $\mathbf{33.7\%}$ & $\mathbf{3.0\%}$ \\
range           & $13.8$--$121.6$\,h & $87.3$--$94.2$\,h \\
STANDBY power   & $4\,570\pm3\,543$\,W & $4\,676\pm343$\,W \\
\bottomrule
\end{tabular}
\end{center}

\textbf{The pipeline is reproducible because of the physics validation,
not despite the clustering.} That is a stronger statement of this
paper's contribution than an unconditional total would be, and it is why
the accept/reject loop is part of the method in
Section~\ref{sec:method:accept_reject} rather than an appendix to it.
Flicker is $0.00\%$ at every seed and the adjusted Rand index against
seed~$0$ stays $\ge0.889$: the \emph{partition} is stable because OFF
and PEAK\_LOAD dominate the record, and it is the identity of the
STANDBY cluster specifically that moves.

\subsubsection{Model-fitting variance is not what makes the economics
uncertain}

Ten seeds of the identical decision configuration (S2, minimum-dwell
labels) give a proposed-policy saving of
$\mathbf{+2.92\pm0.59}$\,USD/yr, $7.8\pm1.6\%$ of oracle, on head-room
of $37.63$\,USD/yr with \emph{no} seed variance at all --- head-room
being a property of the episodes and the coefficients rather than of the
fit. The forecaster's own error is $8\,954\pm156$\,s, about $2.5$\,h on
episode durations, which is why the policy captures under $8\%$ of what
is available.

Set the two sources of uncertainty side by side: model fitting
contributes $\pm0.6$\,USD/yr, day-to-day sampling contributes
$\pm50$\,USD/yr. \textbf{The uncertainty in the economic claim is the
18-day sample, not the model}, and no amount of retraining addresses it.

\subsubsection{State-time totals}

The one-week moving-block bootstrap gives intervals of $\pm15$--$25\%$
of the point estimate on the five-month records
(Table~\ref{tab:multimachine}, per machine). Combined with the $\pm3\%$
conditional labelling spread above, this fixes the precision at which
every state-time number in this paper is quoted: \textbf{two significant
figures, and conditional on the physics battery passing}. Three
significant figures assert a precision neither interval supports.

%  Figures are referenced by bare name; add to the preamble
%      \graphicspath{{../outputs/phase6/figures/}}
%  Bibliography: paper/references.bib
%  Numbers below are generated by  python run_phase5.py  and stored in
%    outputs/phase5/task11/pelletizer-I/forecasting_tests.json
%    outputs/phase5/task11/pelletizer-I/decision_bootstrap_S*.json
%    outputs/phase5/task11/cross_machine_tests.json
%    outputs/phase5/task11/state_block_bootstrap.csv
%    outputs/phase5/task12/pelletizer-I/{forecasting,labelling,decision}/
% =====================================================================

\section{Statistical Analysis}
\label{sec:significance}

Every headline quantity in the preceding sections is a point estimate.
This section reports what happens to those estimates under resampling
and under repeated fitting, and --- for two of the designs used here ---
what they are incapable of detecting in the first place.

\subsection{Four decisions that determine every interval below}
\label{sec:significance:method}

\textbf{The resampling unit is never a row.} A $1$\,Hz industrial record
is massively autocorrelated, and an i.i.d.\ bootstrap over rows would
treat $5.5$ million readings as $5.5$ million independent observations.
Measured on a synthetic AR(1) series with $\rho=0.999$, the i.i.d.\
bootstrap returns a standard error $25\times$ too small. Three
dependence-aware units are used instead, each chosen from the structure
of the quantity being estimated:

\begin{center}
\small
\begin{tabular}{@{}lll@{}}
\toprule
quantity & unit & why \\
\midrule
state-time totals & one-week moving blocks~\cite{kunsch1989jackknife}
  & duty cycle closes over a week \\
decision savings & factory days
  & the restart cap is per calendar day \\
forecast errors & test windows, \emph{paired}
  & one problem solved twice \\
\bottomrule
\end{tabular}
\end{center}

\textbf{A day drawn twice counts as two days.} In the decision bootstrap
the resampled copies are re-keyed. Leaving the label unchanged would let
two copies of one day share a single day's restart budget, silently
tightening the very constraint under evaluation.

\textbf{Every policy is scored on the same resample.} Drawing
independent resamples per policy would break the pairing and inflate the
interval of every contrast --- and the contrasts, not the individual
intervals, are what the decision claim rests on
(Section~\ref{sec:significance:decision}).

\textbf{Statistical and economic significance are separate columns.} A
paired test over $5\,449$ windows detects differences far below anything
that changes a decision, so every comparison carries a $p$-value, an
effect size, and a threshold fixed before the tests were run: $1$\,s of
per-window MAE, $0.02$ of $F_1$, $5$\,USD/yr. A result is allowed to be
\emph{detectable but immaterial}, and two of the twenty-one forecasting
contrasts are. Multiplicity is controlled by Holm's step-down
procedure~\cite{holm1979simple}, which achieves the same family-wise
error rate as Bonferroni and is uniformly more powerful;
Benjamini--Hochberg~\cite{benjamini1995controlling} is reported beside
it for the exploratory families, and both are shown next to the raw
$p$-value so the cost of the correction stays visible.

\subsection{Forecasting}
\label{sec:significance:forecast}

Seven models, $5\,449$ identical held-out windows. Every stored
prediction array was re-scored and matched against its recorded metrics
before any test ran, and all seven sit on a byte-identical vector of
targets, which is what makes the pairing legitimate.

\paragraph{The pre-registered test.}
The comparison named in advance was the proposed Seq2Seq LSTM against
the untrained reference. It resolves decisively, and in the direction
opposite to the one it was written to check
(Table~\ref{tab:sig:headline}). The design was set up to ask whether a
small \emph{positive} gap was real; there is no positive gap. The $F_1$
interval is computed on seed~$0$ of each model, which is the
\emph{best} of the sequence model's five seeds --- and even granting it
its most favourable draw the interval excludes zero in persistence's
favour. On the seed mean the gap is five times wider.

\begin{table}[t]
\centering
\caption{The pre-registered comparison: proposed sequence model versus
the untrained reference, on $5\,449$ paired windows.}
\label{tab:sig:headline}
\small
\begin{tabular}{@{}lrrr@{}}
\toprule
quantity & Seq2Seq & \textsc{persistence} & $p$ \\
\midrule
per-window MAE (s)        & $20.90$ & $12.95$ & $5.7\times10^{-44}$ \\
horizon-step accuracy     & \multicolumn{2}{c}{$-0.0574$ (worse)} & $1.5\times10^{-113}$ \\
STANDBY $F_1$, seed mean  & $0.561$ & $0.681$ & --- \\
STANDBY $F_1$, seed 0     & $0.660$ & $0.681$ & CI $[-0.042,-0.003]$ \\
\bottomrule
\end{tabular}
\end{table}

\paragraph{The family.}
Of the $21$ pairwise contrasts, $19$ are significant after Holm
correction and $17$ of those also clear the $1$\,s practical threshold.
Two conclusions survive both filters. \emph{Persistence beats all five
neural models}, at $p_{\text{Holm}}\le6\times10^{-31}$ and by $4.5$--$8.0$\,s of
MAE. \emph{XGBoost beats persistence}, by $1.45$\,s at
$p_{\text{Holm}}=7.0\times10^{-4}$ with rank-biserial correlation
$0.24$ --- the single contrast in the entire family in which training
beats not training.

\paragraph{Friedman over seeds, and how little it can see.}
Blocking on the five seeds and treating the six trained models as
treatments gives $\chi^2=12.77$, $p=0.026$, Kendall's $W=0.511$
\cite{friedman1937use,demsar2006statistical}. XGBoost takes mean rank
$1.00$ --- it is the best model on every seed block --- and Seq2Seq
takes $5.00$, the worst. But the Nemenyi critical difference is $3.37$
on a $1$--$6$ rank scale, so \textbf{the only pair this diagram
separates is XGBoost against Seq2Seq} ($4.00$ apart) and everything else
is joined (Figure~\ref{fig:cd}). Measured power at this design is $0.14$
for a $1$-sd effect and $0.53$ at $2$\,sd. The seed-blocked comparison
is nearly blind, and the confirmatory conclusions above do not rest on
it: those are paired over $5\,449$ windows, not over $5$ seeds.

\begin{figure}[t]
  \centering
  \includegraphics[width=\linewidth]{fig_p6_cd_diagram}
  \caption{Critical-difference diagram over five seed blocks. Models
  joined by a bar are not separable at $\alpha=0.05$. The bar spans
  almost the whole scale because $n=5$ blocks buys a critical difference
  of $3.37$ ranks.}
  \label{fig:cd}
\end{figure}

\subsection{The decision layer}
\label{sec:significance:decision}

Day-block bootstrap, $2\,000$ draws, on the minimum-dwell labelling:
$394$ held-out episodes over \textbf{18 factory days} spanning $73.6$
calendar days.

\begin{table}[t]
\centering
\caption{Annual saving against the plant's own behaviour (USD/yr per
machine), with $95\%$ day-block bootstrap intervals. Every policy is
scored on the same resample.}
\label{tab:sig:decision}
\small
\setlength{\tabcolsep}{3pt}
\begin{tabular}{@{}lrrr@{}}
\toprule
Policy & S1 ($10$\,min) & S2 ($60$\,min) & S3 ($240$\,min) \\
\midrule
never shut down   & $-219.8$\,{\footnotesize$[-334,-132]$} & $-157.1$\,{\footnotesize$[-276,-65]$} & $+68.6$\,{\footnotesize$[-84,+213]$} \\
immediate         & $-219.7$\,{\footnotesize$[-338,-126]$} & $-307.5$\,{\footnotesize$[-430,-212]$} & $-623.4$\,{\footnotesize$[-774,-478]$} \\
static break-even & $-166.7$\,{\footnotesize$[-290,-71]$}  & $-66.0$\,{\footnotesize$[-196,+24]$}  & $+6.3$\,{\footnotesize$[-170,+168]$} \\
ski-rental        & $-46.7$\,{\footnotesize$[-93,-7]$}     & $+0.0$\,{\footnotesize$[-25,+29]$}    & $+30.4$\,{\footnotesize$[-79,+155]$} \\
\textbf{proposed} & $\mathbf{-28.5}$\,{\footnotesize$[-79,+1]$} & $\mathbf{+2.9}$\,{\footnotesize$[-56,+46]$} & $\mathbf{+94.5}$\,{\footnotesize$[-58,+237]$} \\
oracle            & $+25.5$\,{\footnotesize$[+15,+38]$}    & $+37.6$\,{\footnotesize$[+17,+61]$}   & $+160.8$\,{\footnotesize$[+65,+266]$} \\
\bottomrule
\end{tabular}
\end{table}

Table~\ref{tab:sig:decision} has to be read twice, and the two readings
give different answers.

\textbf{Read row by row, no policy but the oracle demonstrates a
saving.} Every other interval against the status quo spans zero,
including the proposed policy's headline $+2.9$ USD/yr inside
$[-55.9,+46.2]$. The point estimate the earlier sections quote is
indistinguishable from zero, and always was; the bootstrap supplies the
interval that makes it visible.

\textbf{Read as paired contrasts, one result survives}
(Table~\ref{tab:sig:contrasts}). Because every policy is scored on the
same resample, a difference between two rows is far better determined
than either row: the pairing removes the day-to-day variance that makes
the individual intervals wide. The proposed policy beats the static
break-even rule it was built to replace by $+68.9$ USD/yr
($p=0.001$, interval excluding zero, above the $5$ USD/yr economic
threshold) in S2 and by $+138$ and $+88$ USD/yr in S1 and S3. It does
\emph{not} beat forecast-free ski-rental, and it cannot be shown to beat
the plant's own behaviour.

\begin{table}[t]
\centering
\caption{Paired contrasts, same resample. The claim the decision layer
supports is the third row, and only the third row.}
\label{tab:sig:contrasts}
\small
\begin{tabular}{@{}lrrr@{}}
\toprule
Contrast & S1 & S2 & S3 \\
\midrule
proposed $-$ status quo   & $-28.5$ & $+2.9$ & $+94.5$ \\
                          & \footnotesize$p=0.072$ & \footnotesize$p=0.828$ & \footnotesize$p=0.196$ \\
proposed $-$ ski-rental   & $+18.2$ & $+2.8$ & $+64.1$ \\
                          & \footnotesize$p=0.553$ & \footnotesize$p=0.729$ & \footnotesize$p=0.212$ \\
\textbf{proposed $-$ static} & $\mathbf{+138.2}$ & $\mathbf{+68.9}$ & $\mathbf{+88.3}$ \\
                          & \footnotesize$\mathbf{p<0.001}$ & \footnotesize$\mathbf{p=0.001}$ & \footnotesize$\mathbf{p<0.001}$ \\
proposed $-$ oracle       & $-54.0$ & $-34.8$ & $-66.3$ \\
                          & \footnotesize$p<0.001$ & \footnotesize$p<0.001$ & \footnotesize$p<0.001$ \\
\bottomrule
\end{tabular}
\end{table}

\textbf{The oracle gap is also real}, in all three scenarios: the
head-room is genuine and mostly uncaptured --- the proposed policy takes
$7.6\%$ of it in S2.

\textbf{The width is the sample, not the estimator.} Eighteen factory
days is the binding constraint: the day-block bootstrap has eighteen
units to resample, and no refinement of the estimator narrows these
intervals. Any claim of economic benefit from this layer needs a longer
record, and this paper does not make one.
Figure~\ref{fig:decision-forest} shows all three scenarios at once,
with the intervals that cross zero drawn so that they cannot be
mistaken for the ones that do not.

\begin{figure}[t]
  \centering
  \includegraphics[width=\linewidth]{fig_p6_decision_forest}
  \caption{Decision-layer intervals in the three scenarios. Intervals
  crossing zero are drawn grey and hollow; only the oracle's excludes
  zero in every scenario.}
  \label{fig:decision-forest}
\end{figure}

\subsection{Cross-machine tests, and a design that cannot detect
anything}
\label{sec:significance:cross}

Three paired comparisons across the eight machines are reported in
Table~\ref{tab:sig:cross}, and the third column of that table is as
important as the second.

\begin{table}[t]
\centering
\caption{Paired signed-rank tests across machines. The $p$-value floor
is the smallest two-sided value the test can return at that $n$.}
\label{tab:sig:cross}
\small
\setlength{\tabcolsep}{4pt}
\begin{tabular}{@{}lrrrrr@{}}
\toprule
Comparison & $n$ & $\Delta F_1$ & $p$ & floor & power \\
\midrule
own vs \textsc{persistence}, all      & 8 & $+0.194$ & $0.016$ & $0.008$ & $0.41$ \\
own vs \textsc{persistence}, motors   & 4 & $+0.048$ & $0.250$ & $\mathbf{0.125}$ & $\mathbf{0.00}$ \\
local vs transferred                  & 8 & $+0.384$ & $\mathbf{0.008}$ & $0.008$ & $d=1.75$ \\
\bottomrule
\end{tabular}
\end{table}

\textbf{Training beats not training across the eight machines, but the
effect lives in the auxiliary plant.} The $+0.194$ mean difference is
dominated by exhaust-fan~I ($0.982$ against $0.288$) and dpc~I ($0.402$
against $0.012$), machines whose ``STANDBY'' cluster is really their off
state and is therefore trivially predictable once a model has seen it
(Section~\ref{sec:multimachine}). On the four motor-driven machines, the
ones on which the state model is physically valid, the difference is
$+0.048$ and cannot be tested.

\textbf{``Cannot be tested'' is meant literally.} With four paired
observations the smallest two-sided $p$-value the signed-rank test can
return is $2/2^{4}=0.125$, which is above $\alpha=0.05$. No data, and no
effect size however large, can produce a significant result at $n=4$.
\textbf{A non-significant verdict here is a statement about the design,
not about the models}, and it is reported in exactly those terms rather
than as ``no significant difference''.
Figure~\ref{fig:power} puts numbers on both designs: at $n=8$ the paired
test reaches $80\%$ power only above $\approx1.4$\,sd; at $n=4$ it never
does, at any effect size.

\textbf{The transfer degradation is real.} Locally fitted sequence
models beat transferred ones by $0.384$ of $F_1$ with $p$ at the floor
and $d=+1.75$. This is the one cross-machine claim of
Section~\ref{sec:crossmachine} that survives testing outright.

\begin{figure}[t]
  \centering
  \includegraphics[width=\linewidth]{fig_p6_power_curve}
  \caption{Power of the paired signed-rank test against effect size for
  $n=4$ and $n=8$. The marker is the effect actually observed on the
  four motor-driven machines ($d=0.56$), whose measured power is $0.00$
  --- readable off the axis rather than taken on trust.}
  \label{fig:power}
\end{figure}

\subsection{Repeated fitting, and the result that matters most}
\label{sec:significance:seeds}

\subsubsection{The labelling is not reproducible from the seed alone}

Re-fitting the state model at ten seeds on one deterministic
preprocessing of pelletizer~I gives a STANDBY total of
$84.6\pm28.5$\,h --- a coefficient of variation of $\mathbf{33.7\%}$ on
the headline state-time quantity of the paper. Three seeds in ten
converge on a cluster that is not standby at all:

\begin{center}
\small
\begin{tabular}{@{}rrrrl@{}}
\toprule
seed & physics & STANDBY & mean power & what it found \\
\midrule
6 & $0.80$ & $121.6$\,h & $12\,961$\,W & a loaded running state \\
8 & $0.60$ & $63.2$\,h  & $0.8$\,W     & the OFF state \\
9 & $0.60$ & $13.8$\,h  & $1.1$\,W     & the OFF state \\
\bottomrule
\end{tabular}
\end{center}

A ``standby'' cluster at $0.8$\,W drawing $0.06\%$ of rated current is
the machine switched off; one at $13$\,kW is the machine working. EM
initialisation decides which, and in $30\%$ of initialisations it
decides wrongly.

\textbf{What rescues the result is the physics battery, and this is the
constructive finding.} The tests of
Section~\ref{sec:method:accept_reject} score exactly those three seeds
at $0.60$--$0.80$ and the other seven at $1.00$, so the accept/reject
decision is exact on this record. Conditioning on a passing score
collapses the spread:

\begin{center}
\small
\begin{tabular}{@{}lrr@{}}
\toprule
 & all $10$ seeds & the $7$ that pass \\
\midrule
STANDBY hours   & $84.6\pm28.5$\,h & $\mathbf{92.4\pm2.8}$\,h \\
coefficient of variation & $\mathbf{33.7\%}$ & $\mathbf{3.0\%}$ \\
range           & $13.8$--$121.6$\,h & $87.3$--$94.2$\,h \\
STANDBY power   & $4\,570\pm3\,543$\,W & $4\,676\pm343$\,W \\
\bottomrule
\end{tabular}
\end{center}

\textbf{The pipeline is reproducible because of the physics validation,
not despite the clustering.} That is a stronger statement of this
paper's contribution than an unconditional total would be, and it is why
the accept/reject loop is part of the method in
Section~\ref{sec:method:accept_reject} rather than an appendix to it.
Flicker is $0.00\%$ at every seed and the adjusted Rand index against
seed~$0$ stays $\ge0.889$: the \emph{partition} is stable because OFF
and PEAK\_LOAD dominate the record, and it is the identity of the
STANDBY cluster specifically that moves.

\subsubsection{Model-fitting variance is not what makes the economics
uncertain}

Ten seeds of the identical decision configuration (S2, minimum-dwell
labels) give a proposed-policy saving of
$\mathbf{+2.92\pm0.59}$\,USD/yr, $7.8\pm1.6\%$ of oracle, on head-room
of $37.63$\,USD/yr with \emph{no} seed variance at all --- head-room
being a property of the episodes and the coefficients rather than of the
fit. The forecaster's own error is $8\,954\pm156$\,s, about $2.5$\,h on
episode durations, which is why the policy captures under $8\%$ of what
is available.

Set the two sources of uncertainty side by side: model fitting
contributes $\pm0.6$\,USD/yr, day-to-day sampling contributes
$\pm50$\,USD/yr. \textbf{The uncertainty in the economic claim is the
18-day sample, not the model}, and no amount of retraining addresses it.

\subsubsection{State-time totals}

The one-week moving-block bootstrap gives intervals of $\pm15$--$25\%$
of the point estimate on the five-month records
(Table~\ref{tab:multimachine}, per machine). Combined with the $\pm3\%$
conditional labelling spread above, this fixes the precision at which
every state-time number in this paper is quoted: \textbf{two significant
figures, and conditional on the physics battery passing}. Three
significant figures assert a precision neither interval supports.
